\section{Introduction}

Underwater imagery is essential for marine observation, environmental monitoring, offshore inspection, and autonomous robotics. However, the optical properties of water introduce severe degradation: red wavelengths attenuate rapidly, blue and green channels dominate the scene, scattering reduces contrast, and suspended particles produce haze-like veiling. These factors compromise visibility and can significantly degrade downstream tasks such as object detection, segmentation, and visual tracking. Unlike terrestrial imaging, underwater image enhancement must jointly address spatially varying illumination, wavelength-dependent absorption, and color distortion.

Traditional restoration methods commonly rely on handcrafted priors, physical imaging models, and transmission estimation. Representative approaches such as dark-channel priors, dehazing formulations, and color-balance strategies have been effective in limited settings but often fail when the underwater scene violates assumptions of homogeneous medium properties or stable illumination. In addition, these methods can be sensitive to noise, require manual parameter tuning, and are computationally expensive in practical deployment.

Recent years have witnessed a transition from physics-driven models toward data-driven deep learning solutions. Early methods explored generative adversarial networks and end-to-end image-to-image translation to learn restoration mappings directly from paired or unpaired underwater images. As the field matured, more specialized architectures were proposed to model haze, contrast loss, and color cast jointly. Nevertheless, many existing networks remain limited because they either focus primarily on restoration or primarily on color correction rather than performing both tasks in a unified pipeline. This motivates the development of a dual-branch design that explicitly separates the restoration and color-correction objectives while enforcing a shared global representation.

The main contribution of this work is a dual-branch deep learning framework for underwater image restoration and color correction. The first branch is responsible for reconstructing scene radiance and suppressing scattering effects by estimating transmission-aware residuals. The second branch learns channel-wise color compensation to neutralize blue-green dominance and recover natural chroma. A fusion module then combines both outputs using adaptive weighting, producing a final enhanced image with improved contrast, sharper structure, and balanced color distribution. We evaluate the proposed method on the UIEB benchmark and the challenging underwater subset, and results show consistent gains in both objective metrics and perceptual quality compared with strong baselines.

The remainder of the paper is organized as follows. Section~\ref{sec:related_work} reviews representative methods in underwater image enhancement and illustrates the field's transition from physical models to deep learning. Section~\ref{sec:methodology} presents the formulation, architecture, and training objective. Section~\ref{sec:experiments} reports the experimental setup, benchmark datasets, evaluation metrics, and comparative results. Finally, the paper concludes with a discussion of the implications and future research directions.
